\section{Finite registers and the numerical BP2 frontend}
\label{sec:source-registers}

The source reduction is an executable composition, with no supplied bound on
source running time.  Its intermediate mechanisms follow ais523's
Amnesiac From Minsk~\cite{amnesiac},
Waterfall Model~\cite{waterfall},
and Brainpocalypse II~\cite{bp2}
constructions; the Waterfall account also credits zzo38's independent model.
The explicit register-frontend tables, compiler composition, and proofs
described here are this development's contributions.  The generic CTS-to-S
encoding and root-reset controller are supplied by
\href{https://github.com/cstrawberry/predictive-universe/tree/85a867988442fc423279341200f81634a1e65582/docs/paper/related/pure_s_universality}{Cinematic Strawberry's pinned construction}.

A source machine has $R\geq1$ natural registers, $Q\geq1$ labels, an entry
label, and instructions $\operatorname{INC}(i,t)$,
$\operatorname{DECJZ}(i,t_+,t_0)$, and $\operatorname{HALT}$.
Zero testing leaves a zero register unchanged.  BP2 instead executes a finite
list of $+i,-i$: positive decrements advance normally, while a zero decrement
sets that counter to one and restarts at command zero, retaining all other
counters.  Falling off the list halts; all BP2 counters initially equal zero.

\begin{theorem}[Exact register frontend]
For every such $M$ and input tuple $a$, the total function
\code{SOnlyMachineNat.compile} produces a nonempty finite BP2 program $B(M,a)$
with dense counter labels.  It halts exactly when $M(a)$ halts.  For each
$v\in\NN$, it halts with first counter $2v+3$ exactly when $M(a)$ halts with
register zero equal to $v$.
\end{theorem}

\begin{proof}[Construction and proof]
First use positive amnesiac counters: $+i$ increments and follows a fixed
successor $p_i$; $-i$ decrements above one and follows $s_i$, or leaves one
unchanged and follows $f_i$.  Lean stores these values minus one.  Keep
$D_i=r_i+1$ and allocate five helpers $I_q,S_q,F_q,U_q,V_q$ at \emph{every}
label, initially one.  Let $E(q)$ be $+I_q$, $+S_q$, or halt according to the
source instruction.  The successor table is
\[
\begin{array}{c|ccc}
 &p&s&f\\\hline
D_i&-I_0&-S_0&-F_0\\
I_q&+D_i&E(t)&-I_{q+1}\\
S_q&+F_q&-V_q&-S_{q+1}\\
F_q&-D_i&-U_q&-F_{q+1}\\
U_q&-F_0&E(t_+)&+V_q\\
V_q&-S_0&E(t_0)&+U_q
\end{array}
\]
Here instruction-dependent entries apply to the appropriate instruction;
inapplicable instruction-dependent entries and end-of-chain fallbacks are halt.  Scans visit all
labels.  A chain with exactly one marked value two passes the ones, restores
the mark to one, and takes precisely its success edge.  An increment marks
$I_q$, increments $D_i$, and scans $I$.  A decrement marks $S_q,F_q$, tests
$D_i$, then uses the $S,V,U,F$ or symmetric $F,U,V,S$ route to restore every
helper and select $t_+$ or $t_0$.  Thus the instruction-boundary invariant is
exact, including shared data registers and repeated call sites
(\code{SOnlyMachine.macro_simulation}).

For $k=R+5Q$ amnesiac counters, introduce $4k+1$ Waterfall clocks: data $X_i$,
action clocks $Q_a$ for $a=+i,-i,H$, and recovery clocks $T_i$.
At normalized action $(c,a)$ set $X_i=2c_i$, $Q_a=0$, and every other control
clock to four.  With $e_Z$ a unit vector, the signed trigger rows are
\[
\begin{aligned}
b_{Q_{+i}}&=2e_{X_i}+4e_{Q_{+i}}-4e_{Q_{p_i}},&
b_{Q_{-i}}&=-2e_{X_i}+4e_{Q_{-i}}-3e_{T_i},\\
b_{X_i}&=2e_{X_i}+3e_{T_i}-4e_{Q_{f_i}},&
b_{T_i}&=\mathbf1+3e_{T_i}-4e_{Q_{s_i}}.
\end{aligned}
\]
Use $A_Z=b_Z+5\cdot\mathbf1$ and $A_{Q_H}=0$, adding coincident terms.
All nonhalt entries lie in $[1,9]$.  Subtracting the common minimum after a
trigger cancels its uniform offset.  An increment reaches the next normalized
action; a decrement makes $X_i=2c_i-2$ and $T_i=1$.  If $c_i=1$, $X_i$ is
uniquely minimal and restores the failure state; otherwise $T_i$ is uniquely
minimal and its recovery row restores the successful-decrement state.
These arithmetic identities prove tie-free simulation
(\code{SOnlyMachineClock.step_simulation}).

BP2 implements each clock with a value cell $v_i$ and pending marker $z_i$,
plus one marker $m$.  For a nonnegative vector $b$, write $\operatorname{Add}(b)$
and $\operatorname{Sub}(b)$ for literal unary increments and decrements.
The initializer is
\[
 \operatorname{Add}(a);-m;+m;\operatorname{Sub}(a),
 \qquad a(v_i)=w_i,\quad a(z_i)=1,\quad a(m)=0.
\]
Its first zero guard preserves initialization across a restart; subsequent
visits cancel exactly, even when working cells are zero.
Trigger $i$ uses the same pattern guarded by $z_i$, adding
$A_i-e_i$ to value cells.  A sweep subtracts one from every $v_i$; tests
$-z_i;-v_i;+v_i;+z_i$ mark the unique zero and restart.  The selected trigger
then compensates for the test's automatic increment and applies exactly $A_i$.
The terminal pair $-m;-m$ restarts ordinary sweeps.  The halt trigger instead
adds $2\cdot\mathbf1-e_h$ and increments $m$; one final sweep falls off with
$v=y+\mathbf1$, every $z_i=1$, and $m=0$.  Hence the first value is
$2(r_0+1)+1=2r_0+3$.
Finite positive macro durations and absorbing halt states give reflection at
arbitrary target times, not only checkpoints.  Dense renumbering preserves
the output coordinate.  These are
\code{SOnlyMachineNat.halt_iff} and
\code{SOnlyMachineCorrectness.output_iff}.
\end{proof}

The formal compiler initializes clocks at $V+5$, whereas the earlier selective
Python construction uses $V+1$ and allocates helpers only where needed.
The formal count is $c=8k+3$, $k=R+5Q$; its checked bound, with
$N_c=4k+1$ and $s=\sum_i a_i$, is
\[
 L\leq2c(2s+9)+N_c(18c+2)+c+4N_c+4
 \quad\text{(\code{SOnlyMachineBounds.compile_length_bound}).}
\]
The two variants preserve semantics but are not byte-for-byte identical.

\section{Programmed UT19 execution and its first event}
\label{sec:source-ut19}

The next reduction uses ais523's CC0
UT19 construction~\cite{ut19}.
Its literal productions, in published one-based notation, are
\[
\begin{array}{r@{\,\mapsto\,}l r@{\,\mapsto\,}l r@{\,\mapsto\,}l}
1&(2,3)&2&(4,4)&3&(18,4)\\
4&(1,1,19)&5&(7,9)&6&(8,9)\\
7&(10,10)&8&(11,10)&9&(18,10)\\
10&(5,6)&11&(6,5,19)&12&(14,14,14,14)\\
13&(15)&14&(16,16)&15&(16,17)\\
16&(12,13)&17&(12,12,12,12)&18&(18)\\
19&\epsilon&\multicolumn{2}{c}{}&\multicolumn{2}{c}{}
\end{array}
\]
Word powers denote literal repetition. A microstep removes one symbol, appends its production only in \emph{take}
phase, and toggles take/skip.  The designated event is a \emph{prestate} with
take phase and head 18; a skipped 18 is not an event.

For a BP2 program of length $L$ and $c$ counters, add a halt counter $H=c+1$.
Choose the least power of two $N\geq L+3$ and put $W=2N$.
The seed is
\[
19\,M_1K_1\cdots M_cK_cM_HK_HM_{H+1},
\qquad K_i=(12,13)^2.
\]
Each memory has $W$ cells.  The expanded program is dummy, the $L$ commands,
halt-A, halt-B, and padding to $N$ slots.  Every slot uses two halfcommands;
each halfcommand comprises Command, Parity, and Reset passes, where a pass
consumes exactly its entrance word.

\begin{lemma}[Finite programmed-memory window]
For any prescribed $W$ width parities there is an explicit initializer giving
those observations at times $0,\ldots,W-1$, independently of the incoming
Command alignment history.  Odd Reset restores this initializer exactly.
\end{lemma}

\begin{proof}[Proof sketch]
A cell with dynamic bit $h_i$ is $(5,6)$ or $(6,5,19)$; a fixed inverter bit
$a_i=1$ adds $(1,1,19)$.  Put $c_i=h_i\oplus a_i$.
Direct production calculation gives, under normal Parity and Reset,
$c'=Pc\oplus g\mathbf1$, where $P$ is inclusive prefix XOR and $g=1$ means
skip alignment.  Width is observed \emph{before} updating:
$v_t=\bigoplus_i c_{t,i}$.
For $W=2^d$, the unforced observation transform is
\[
 v_t=\bigoplus_{i:\,(i\mathbin{\&}t)=t}a_i,
 \qquad
 a_i=\bigoplus_{j:\,(j\mathbin{\&}i)=i}v_j.
\]
The superset transform is its own inverse over $\mathbb F_2$.
The formal proof uses even/odd index decomposition and identifies its recursive
transform with the increasing-stride butterfly.  The extensionally equal increasing-stride implementation uses
$d=\log_2W$ stages and $W\log_2W/2$ bit XORs. This count concerns that
implementation; the output-equivalence theorem is not an exact operation-count
theorem for direct evaluation of the recursive Lean definition.  An alignment disturbance in
update $t$ first affects observation $t+W$; reset sets every $h_i$ to zero.
These facts are \code{SOnlyTemporalMemory.forced_initializer_correct},
\code{initializer_eq_increasing_strides}, and
\code{SOnlyMemory.memory_reset_vector}.  The emitted choice
$d=\lfloor\log_2(L+2)\rfloor+2$ satisfies
$2L+3<W\leq4L+8$ by \code{SOnlySimulation.compiler_capacity} and
\code{compiler_width_bound}.
\end{proof}

Memory widths are adjacent XOR differences of the desired component
alignments, so prefix XOR telescopes to exactly those alignments.
\code{SOnlySimulation.scheduledWidth_eq_published} checks equality with the
literal width recipes, including unused halt-B and padding slots.
For $\mathcal C(n)=(12,13)^{2^n}$, ordinary BP2 value $x$ is
$\mathcal C(2x+1)$.  A halfcommand increments $n$ when Command starts take;
a positive-exponent decrement subtracts one under even Reset.
Thus $+i$ increments all exponents, then increments $i$ and decrements the
others.  Similarly, $-i$ decrements all, then decrements only $i$ and
increments the others.  These implement changes of two in the target
exponent and zero elsewhere.

The zero branch requires more than ordinary decrement arithmetic.
Decrementing $\mathcal C(0)$ produces a singleton 15 in Parity.  At a source
zero decrement it is selected, creates odd Reset, and leaves the protected
word $D=(12,12,12,12)$.  Every other counter is incrementing in that
halfcommand: no positive decrement is exposed to odd Reset, which would
otherwise produce the wrong protected form.  All memories reset exactly;
any 18 they emit is skipped.  The restarted dummy decrements all and then
increments all, but $D$ responds to the first half by becoming
$\mathcal C(2)$, then $\mathcal C(3)$.  The target therefore becomes one,
as BP2 requires, while other values are retained.  The initial leading 19
supplies the same skip alignment for the first dummy.  This proves closure
through arbitrarily many restarts
(\code{SOnlySimulation.source_zero_decrement_run}).

At command $i$, memory time is $2i+2$.  Each epoch either resets or reaches
halt-A's last needed observation $2L+3<W$, so the memory lemma is never used
past its protected window.  Halt-A decrements all exponents, then increments
ordinary counters and decrements the halt counter with odd Parity.  The latter's
singleton 15 is skipped and disappears; Reset begins take with head 18.
Exactly five passes after halt-A entrance, the event queue is
\[
 R_1\,16^{4^{x_1+1}}\cdots R_c\,16^{4^{x_c+1}}R_HT,
\]
where all separators are nonempty and contain no 16.  The $x_i$ are the final
BP2 values.  Thus maximal 16-runs give precisely the retained counters
(\code{SOnlySimulation.finalEvent_readout}); repeated division by four
recovers them.  The first counter then yields $(x_1-3)/2$ for register output
(\code{finalEvent_readMachineValue}).

Each preceding simulated step is a finite nonempty, event-free run.
Unconsumed entrance symbols prevent exhaustion inside a pass, and memories
supply nonempty output at its end.  Positive simulation durations make
nonhalting checkpoints cofinal, excluding an event at \emph{every} finite
microstep.  Consequently \code{SOnlySimulation.source_halting_iff_event}
is a two-way unbounded theorem, and \code{first_event_exact} identifies the
actual first event queue, rather than only an expected macro-boundary.

\section{The fixed 38-phase interface and conventional tapes}
\label{sec:source-cts}

Encode symbol $s$ by $\eta(s)=0^{s-1}1\,0^{19-s}$.
The fixed CTS has the nineteen encoded productions followed by nineteen empty
appendants.  Take and skip boundaries are phases zero and nineteen.

\begin{lemma}[Exact block and event transfer]
Nineteen CTS steps simulate one UT19 microstep.  At every CTS time $t$, the
phase-17/head-one event holds exactly when UT19 is selected at microstep
$\lfloor t/19\rfloor$ and $t\bmod19=17$.  A first UT19 event at $n$ therefore
becomes the first CTS event at $19n+17$.
\end{lemma}

\begin{proof}[Proof sketch]
After consuming a prefix, FIFO order leaves its unconsumed suffix ahead of
all generated appendants.  Scanning one 19-bit block appends exactly its
production in the take half, or nothing in the skip half, and toggles the
boundary phase.  The remaining original bits prevent premature exhaustion.
The only one in $\eta(s)$ is at offset $s-1$; phase 17 can therefore signal only
selected symbol 18.  Division with remainder covers arbitrary times.
These are \code{SOnlySource.trajectory_simulation},
\code{event_at_arbitrary_time}, and \code{first_event_exact}.
\end{proof}

At the event, the current CTS word is $10\,\eta(\text{tail})$, not a block-aligned
encoding.  Restoring the known seventeen-zero prefix reconstructs the whole
UT19 event word; \code{SOnlyResultBits.decodeEventWord_correct} proves this
structural readout.  Composition gives
\code{SOnlySimulation.compiled_halting_iff_cts_event} and
\code{compiled_noPrematureEmpty}, the exact source interfaces consumed by the
fixed S backend.  In particular, UT19's rule $18\mapsto18$ is an event marker,
not a claim that the tag queue or S term normalizes.

Finally, the source need not be assumed universal by convention.
\code{SOnlyTapeStacks} treats a deterministic finite-alphabet machine on an
integer-indexed tape, with left, stay, and right moves and absent transition
meaning halt.  Finite control stores the state and scanned symbol; nearest-first
left and right stacks represent every remaining tape cell, with empty pop
returning blank.  For alphabet size $A$, use base $B=A+1$ and
\[
 \operatorname{code}([])=0,\qquad
 \operatorname{code}(a::u)=(a+1)+B\operatorname{code}(u).
\]
Nonzero digits preserve explicit blanks and make this code injective.
Literal INC/DECJZ macros implement push by multiplication and addition, and
pop by a finite-control residue cycle and quotient transfer.  Registers zero
and one hold the stacks; register two is scratch, restored to zero.
The assembled table has
$|Q_{\rm TM}|A(5(A+1)+5)+1$ labels and depends only on the machine; its entry
selects the input's first symbol and the right register contains the remaining
input.  Positive-duration macro execution and cofinal checkpoints reflect
halting even inside target macros
(\code{SOnlyTapeCompiler.input_halting_iff}).

Hence \code{SOnlyTuringUniversality.halting_iff_event} covers arbitrary finite
inputs without a tape-window or runtime premise.  Arbitrary register programs
retain their exact natural output.  For this TM bridge the output is specifically
the finite left-stack code, including retained blanks and visited extent:
\code{left_stack_result_iff_first_event}.  It is a history-dependent observable
of this simulation, not a canonical final tape-output word.
